\documentclass{article}

\usepackage{booktabs}
\usepackage{tabularx}
\usepackage{xltabular}
\usepackage{array}

\title{Development Plan\\\progname}

\author{\authname}

\date{}

\input{../Comments.text}
\input{../Common.text}

\begin{document}

\maketitle

\begin{table}[hp]
\caption{Revision History} \label{TblRevisionHistory}
\begin{tabularx}{\textwidth}{llX}
\toprule
\textbf{Date} & \textbf{Developer(s)} & \textbf{Change}\\
\midrule
Date1 & Name(s) & Description of changes\\
Date2 & Name(s) & Description of changes\\
... & ... & ...\\
\bottomrule
\end{tabularx}
\end{table}

\newpage{}

\wss{Put your introductory blurb here.  Often the blurb is a brief roadmap of
what is contained in the report.}

\wss{Additional information on the development plan can be found in the
\href{https://gitlab.cas.mcmaster.ca/courses/capstone/-/blob/main/Lectures/L02b_POCAndDevPlan/POCAndDevPlan.pdf?ref_type=heads}
{lecture slides}.}

\section{Confidential Information?}

\wss{State whether your project has confidential information from industry, or
not.  If there is confidential information, point to the agreement you have in
place.}

\wss{For most teams this section will just state that there is no confidential
information to protect.}
\section{IP to Protect}

\wss{State whether there is IP to protect.  If there is, point to the agreement.
All students who are working on a project that requires an IP agreement are also
required to sign the ``Intellectual Property Guide Acknowledgement.''}

\section{Copyright License}

\wss{What copyright license is your team adopting.  Point to the license in your
repo.}

\section{Team Meeting Plan}

\wss{This section is for meetings between the team members, and meetings between 
the team and their supervisor/advisor (as appropriate).}

\wss{How often will the team meet? where? when?}

\wss{If the meeting is a physical location (not virtual), out of an abundance of
caution for safety reasons you shouldn't put the location online}

\wss{To help schedule meetings, you can use the lecture/tutorial time when we do 
not have anything else scheduled.}

\wss{How often will you meet with your supervisor/advisor?  when?  where?}

\wss{Will meetings be virtual?  At least some meetings should likely be
in-person.}

\wss{How will the meetings be structured?  There should be a chair for all meetings.  
There should be an agenda for all meetings.}

\section{Team Communication Plan}

\wss{Issues on GitHub should be part of your communication plan.}

\section{Team Member Roles}

\wss{You should identify the types of roles you anticipate, like notetaker,
leader, meeting chair, reviewer.  Assigning specific people to those roles is
not necessary at this stage.  In a student team the role of the individuals will
likely change throughout the year.}

\section{Workflow Plan}

\begin{itemize}
	\item How will you be using git, including branches, pull request, etc.?
	\item How will you be managing issues, including template issues, issue
	classification, etc.?
  \item Use of CI/CD
\end{itemize}

\section{Project Decomposition and Scheduling}

For aiding in documenting and tracing our project deliverables, we will be making use of two GitHub projects that will handle our code and our documentation deliverables. They are linked below: 

\begin{itemize}
    \item \textbf{Code Project:} \url{https://github.com/users/arjunkarthik123/projects/4/views/1}
    \item \textbf{Documentation Project:} \url{https://github.com/users/arjunkarthik123/projects/3}
\end{itemize}

Both projects are split into four sections: Backlog, In Progress, In Review, and Done. As the project progresses, we will come up with items that need to be addressed for both our documentation and later on our codebase. Whenever any team member thinks of an item that needs to be worked on, it will be added to the backlog of the respective projects. Depending on what the team decides, someone will be assigned the new issue, and whenever the person works on it, it will be under the In Progress section. In order to complete the deliverable, it first needs to be reviewed, hence the In Review section. For either the code or documentation, a pull request will be created and linked to the respective issue being addressed, and only when it gets approved will the issue move to the Done section. 

If the issues seem to be larger than anticipated, we will split up the issue using sub-issues and have them assigned accordingly. This is to ensure no single issue is too big to handle by one person alone. 

Table~\ref{TblDeliverables} provides a high-level overview of the key deliverables throughout the academic year (please note that these items were taken directly from the course schedule in the course outline): 

\begin{xltabular}{\textwidth}{lll>{\raggedright\arraybackslash}X}
\caption{Key Deliverables Schedule} \label{TblDeliverables}\\
\toprule
\textbf{Term} & \textbf{Week} & \textbf{Date} & \textbf{Tasks / Deliverables}\\
\midrule
\endfirsthead

\caption[]{Key Deliverables Schedule (Continued)}\\
\toprule
\textbf{Term} & \textbf{Week} & \textbf{Date} & \textbf{Tasks / Deliverables}\\
\midrule
\endhead

\midrule
\multicolumn{4}{r}{\emph{Continued on next page}}\\
\bottomrule
\endfoot

\bottomrule
\endlastfoot

Term 1 & Week 2  & Week of Sep 14, 2026 & \textbf{Deliverable:} Deadline for teams formation, projects selection (end of week)\\
\addlinespace
Term 1 & Week 3  & Week of Sep 21, 2026 & \textbf{Deliverable:} Biosketches due to Avenue dropbox \newline \textbf{TA Meeting:} Review of prob state, PoC plan, dev plan and team charter\\
\addlinespace
Term 1 & Week 4  & Week of Sep 28, 2026 & \textbf{Deliverable:} Problem Statement, Proof of Concept Plan, Development Plan, Team Charter Rev 0 \newline \textbf{Peer Review:} Prob State + PoC plan + Dev Plan + Team Charter peer review\\
\addlinespace
Term 1 & Week 5  & Week of Oct 05, 2026 & \textbf{TA Meeting:} Review SRS+Haz Analysis\\
\addlinespace
Term 1 & Week 6  & Week of Oct 19, 2026 & \textbf{Deliverable:} SRS + Haz Analysis Rev 0 \newline \textbf{Peer Review:} SRS + Haz Analysis peer review\\
\addlinespace
Term 1 & Week 7  & Week of Oct 26, 2026 & \textbf{TA Meeting:} Review VnV Plan\\
\addlinespace
Term 1 & Week 8  & Week of Nov 02, 2026 & \textbf{Deliverable:} VnV Plan Rev 0 \newline \textbf{Peer Review:} VnV Plan peer review\\
\addlinespace
Term 1 & Week 9  & Week of Nov 09, 2026 & \textbf{TA Meeting:} Review Des Doc Rev-1 and CD-Fall\\
\addlinespace
Term 1 & Week 10 & Week of Nov 16, 2026 & \textbf{Deliverable:} Des Doc Rev -1 and CD-Fall \newline \textbf{Peer Review:} Des Doc Rev-1 Peer Review\\
\addlinespace
Term 1 & Week 11 & Week of Nov 23, 2026 & \textbf{Presentation/Demo:} Proof of Concept Demos\\
\addlinespace
Term 1 & Week 12 & Week of Nov 30, 2026 & \textbf{Presentation/Demo:} Proof of Concept Demos\\
\addlinespace
Term 1 & Week 13 & Week of Dec 07, 2026 & \textbf{Presentation/Demo:} Code review interviews\\
\midrule
Term 2 & Week 15 & Week of Jan 11, 2027 & \textbf{TA Meeting:} Review of Des Doc\\
\addlinespace
Term 2 & Week 16 & Week of Jan 18, 2027 & \textbf{Deliverable:} Design Document Rev 0 \newline \textbf{Peer Review:} Module Implementation due\\
\addlinespace
Term 2 & Week 17 & Week of Jan 25, 2027 & \textbf{Peer Review:} Des Document peer review (module implementation reflection)\\
\addlinespace
Term 2 & Week 18 & Week of Feb 01, 2027 & \textbf{Presentation/Demo:} Rev 0 Demo\\
\addlinespace
Term 2 & Week 19 & Week of Feb 08, 2027 & \textbf{Presentation/Demo:} Rev 0 Demo\\
\addlinespace
Term 2 & Week 20 & Week of Feb 22, 2027 & \textbf{Presentation/Demo:} Code review interviews\\
\addlinespace
Term 2 & Week 21 & Week of Mar 01, 2027 & \textbf{TA Meeting:} Review VnV Report+CD-Winter\\
\addlinespace
Term 2 & Week 22 & Week of Mar 08, 2027 & \textbf{Deliverable:} VnV Report + CD-Winter \newline \textbf{Peer Review:} VnV Report + CD-Winter peer review\\
\addlinespace
Term 2 & Week 23 & Week of Mar 15, 2027 & \textbf{Presentation/Demo:} Final Demos\\
\addlinespace
Term 2 & Week 25 & Week of Mar 29, 2027 & \textbf{Deliverable:} EXPO poster pdf and video due; Final Documentation (Rev 1)\\
\addlinespace
Term 2 & Week 26 & Week of Apr 05, 2027 & \textbf{Presentation/Demo:} Capstone EXPO\\
\end{xltabular}

\section{Proof of Concept Demonstration Plan}

What is the main risk, or risks, for the success of your project?  What will you
demonstrate during your proof of concept demonstration to convince yourself that
you will be able to overcome this risk?

\section{Expected Technology}

\wss{What programming language or languages do you expect to use?  What external
libraries?  What frameworks?  What technologies.  Are there major components of
the implementation that you expect you will implement, despite the existence of
libraries that provide the required functionality.  For projects with machine
learning, will you use pre-trained models, or be training your own model?  }

\wss{The implementation decisions can, and likely will, change over the course
of the project.  The initial documentation should be written in an abstract way;
it should be agnostic of the implementation choices, unless the implementation
choices are project constraints.  However, recording our initial thoughts on
implementation helps understand the challenge level and feasibility of a
project.  It may also help with early identification of areas where project
members will need to augment their training.}

Topics to discuss include the following:

\begin{itemize}
\item Specific programming language
\item Specific libraries
\item Pre-trained models
\item Specific linter tool (if appropriate)
\item Specific unit testing framework
\item Investigation of code coverage measuring tools
\item Specific plans for Continuous Integration (CI), or an explanation that CI
  is not being done
\item Specific performance measuring tools (like Valgrind), if
  appropriate
\item Tools you will likely be using?
\end{itemize}

\wss{git, GitHub and GitHub projects should be part of your technology.}

\section{Use of AI and LLMs}

\wss{This section is a continuation of the previous section.  The role of LLMs
in your project as important enough to warrant their own section.  
How will your team be using LLMs for code and documentation development?  
What tools will you be using?  How will you verify the output of your LLMs?}

\section{Coding Standard}

\wss{What coding standard will you adopt?}

\newpage{}

\section{Appendix --- Generative AI Use Disclosure}

\wss{ChatGPT (version 5) was used to brainstorm ideas for the template for this disclose section.}

\wss{This section is about the use of AI to write this document, not your
planned use of AI for your project.  You discuss that topic in one of the
sections in the body of this report.}

\subsection{AI Tools Used}

\wss{Give the name of the AI tool(s) used and the version.}

\subsection{How and Where Used}

\wss{Explain what the tool(s) were used for.  Examples include: brainstorming,
explaining a technical concept, reviewing the document, generating or modifying
text, generating or modifying diagrams, identifying errors or omissions or
inconsistencies.}

\wss{For each of the uses, identify which parts of the document were affected.}

\wss{You do not need to report every interaction, but you should disclose the 
uses that materially contributed to your work}

\subsection{Verification of AI-Generated Content}

\wss{\begin{itemize}
    \item Describe how the team checked AI-generated or AI-assisted material
    before including it in the document.
    \item In particular, verify factual claims, technical assertions,
    requirements, calculations, references, and other information for which an
    AI system may produce incorrect or fabricated results.
    \item \textbf{The team is responsible for the accuracy, completeness,
consistency, and appropriateness of the submitted document, regardless of
whether material was generated by a human or an AI system.}
\end{itemize}}

\newpage{}

\section*{Appendix --- Reflection}

\wss{Not required for CAS 741}

\input{../Reflection.text}

\begin{enumerate}
    \item Why is it important to create a development plan prior to starting the
    project?
    \item In your opinion, what are the advantages and disadvantages of using
    CI/CD?
    \item What disagreements did your group have in this deliverable, if any,
    and how did you resolve them?
\end{enumerate}

\newpage{}

\section*{Appendix --- Team Charter}

\wss{borrows from
\href{https://engineering.up.edu/industry_partnerships/files/team-charter.pdf}
{University of Portland Team Charter}}

\subsection*{External Goals}

\wss{What are your team's external goals for this project? These are not the
goals related to the functionality or quality fo the project.  These are the
goals on what the team wishes to achieve with the project.  Potential goals are
to win a prize at the Capstone EXPO, or to have something to talk about in
interviews, or to get an A+, etc.}

\subsection*{Attendance}

\subsubsection*{Expectations}

\wss{What are your team's expectations regarding meeting attendance (being on
time, leaving early, missing meetings, etc.)?}

\subsubsection*{Acceptable Excuse}

\wss{What constitutes an acceptable excuse for missing a meeting or a deadline?
What types of excuses will not be considered acceptable?}

\subsubsection*{In Case of Emergency}

\wss{What process will team members follow if they have an emergency and cannot
attend a team meeting or complete their individual work promised for a team
deliverable?}

\subsection*{Accountability and Teamwork}

\subsubsection*{Quality} 

\wss{What are your team's expectations regarding the quality
of team members' preparation for team meetings and the quality of the
deliverables that members bring to the team?}

\subsubsection*{Attitude}

\wss{What are your team's expectations regarding team members' ideas,
interactions with the team, cooperation, attitudes, and anything else regarding
team member contributions?  Do you want to introduce a code of conduct?  Do you
want a conflict resolution plan?  Can adopt existing codes of conduct.}

\subsubsection*{Stay on Track}

\wss{What methods will be used to keep the team on track? How will your team
ensure that members contribute as expected to the team and that the team
performs as expected? How will your team reward members who do well and manage
members whose performance is below expectations?  What are the consequences for
someone not contributing their fair share?}

\wss{You may wish to use the project management metrics collected for the TA and
instructor for this.}

\wss{You can set target metrics for attendance, commits, etc.  What are the
consequences if someone doesn't hit their targets?  Do they need to bring the
coffee to the next team meeting?  Does the team need to make an appointment with
their TA, or the instructor?  Are there incentives for reaching targets early?}

\subsubsection*{Team Building}

\wss{How will you build team cohesion (fun time, group rituals, etc.)? }

\subsubsection*{Decision Making} 

\wss{How will you make decisions in your group? Consensus?  Vote? How will you
handle disagreements? }

\end{document}